\setchapterpreamble[u]{\margintoc}
\chapter{Limitations}
\labch{options}

This master’s thesis was subject to several limitations that affected different stages of the project. The following section describes the most relevant methodological, computational, and experimental constraints encountered throughout the study.

The first major limitation derives from the novelty of ST technologies, particularly Xenium. Since these technologies have only been available for approximately four to seven years, the amount of comparable literature and established methodologies remains limited. At the computational level, few standardized pipelines have been specifically developed for Xenium and Visium data, making the implementation of analysis workflows challenging for researchers entering the field. Furthermore, many computational approaches remain experimental and have not yet consistently demonstrated superior performance compared with conventional statistical methods. Consequently, additional benchmarking studies are still required. From a biological and methodological perspective, the available literature related to MIBC is also scarce, and best practices are still evolving. As a result, several analytical decisions had to be based on exploratory approaches or by reproducing strategies from studies with \textbf{partially comparable} objectives or datasets.

Another important limitation concerns the design of the gene panel employed in this study. As discussed in Section \refsec{VisiumVSXenium}, the number of detectable genes directly influences the reliability and resolution of cell-type annotation. Although the selected panel was sufficient to apply unsupervised analytical approaches, the number of genes was not adequate to support robust clustering analyses with a high degree of confidence. Consequently, the resulting annotations may contain an inherent degree of uncertainty that cannot be completely eliminated in the absence of additional transcriptomic information. Similarly, the cell-type labels assigned using SingleR cannot be interpreted as definitive confirmations of cellular identity, since no absolute ground-truth reference was available for validation. Instead, the analysis assigned the most probable cell-type label to groups of cells whose expression profiles more closely resembled a given reference population than any alternative cell type.

The scale and complexity of the dataset also constituted a major limitation throughout the project. Even relatively simple analytical and visualization tasks were computationally demanding due to the large volume of data generated by Xenium. Approximately 5,399,355 cells were analyzed, with an average of 377 detected genes per cell, compared with approximately 99,028 Visium spots and an average of 4,097 unique genes per spot. This substantial increase in the number of observations resulted in highly memory-intensive workflows, complicating code development, debugging, and optimization procedures. Although access to a high-performance computing cluster was available, memory allocation remained a recurrent bottleneck. For example, a minimum of 300 GB of RAM was required to compute quality-control metrics for all 15 samples included in the study. Furthermore, biological variability between samples occasionally produced unexpected computational issues, including errors associated with insufficient data distributions for specific statistical operations\sidenote{Errors like \texttt{Error in \textasciigrave cut\_number()\textasciigrave: Insufficient data values to produce 24 bins.}}. In some cases, visualization procedures were unable to render outputs because of the excessive number of cells present in individual samples. These challenges highlight the strong dependency of computational performance not only on algorithmic implementation, but also on the structural and mathematical properties of the biological data itself.

% The annotation code was too memory heavy as with 600G of memory we could only make 4 samples and a half due to the time limit\sidenote{GOT THIS ERROR: slurmstepd: error: *** JOB 3709460 ON c02 CANCELLED AT 2026-05-06T08:12:47 DUE TO TIME LIMIT ***}. But when we changed the partition to have more time then 300G was insuficcient memory. So we came back again to the higMem partition and increased the number of cpus (from 4 to 6 as the máximum was 8 in this partition) so teh work could be made fatser but stil it only increased to 6 complete samples. So we had to run multiple Jobs with the optimal combination found of hyperparameters and execute it for batched of samples. In contrast the neighbourhod step needed more time than memory as the highMem oartition with 600G only got 6 samples due to time limit but hpc partition with 300G got all the samples as it neded 4h to execute, both atemps with 4 cpus per task. 

Additional constraints were imposed by the resource management policies of the computational cluster. To ensure equitable access to shared resources, the cluster infrastructure was divided into partitions\sidenote{Partitions are sections of the cluster dedicated to specific types of executions.} optimized for different computational requirements. The \texttt{hpc} partition allowed unrestricted execution time but imposed a memory limit of 300 GB, whereas the \texttt{highMem} partition provided unrestricted memory allocation but limited job execution to a maximum duration of 24 hours. Several stages of the analysis simultaneously required high memory consumption and long execution times, making the selection of an appropriate partition particularly challenging. For instance, the annotation workflow exceeded the execution-time limit even when allocated 600 GB of RAM within the \texttt{highMem} partition, whereas the \texttt{hpc} partition did not provide sufficient memory resources to complete the same analysis. Increasing the number of CPUs improved performance only partially, and the workflow ultimately required dividing the dataset into multiple batches of samples to ensure successful execution. In contrast, the neighbourhood analysis was more constrained by execution time than by memory usage and could therefore be completed within the \texttt{hpc} partition. Future methodological improvements could involve implementing more efficient parallelization strategies. Given the structure of the workflow and the computational resources available within the cluster, distributing the processing of independent samples across multiple CPU groups could potentially reduce execution times substantially.


Reproducibility represented an additional methodological limitation. Functions that depend on random number generation (RNG) may produce different outputs across independent executions unless a fixed random seed is explicitly defined using \texttt{set.seed()}. Although this issue was addressed in the neighbourhood and Cohen’s D analyses, it was not initially implemented in the annotation workflow. The reproducibility limitation originated from the use of \texttt{SingleR} with the parameter \texttt{aggr.ref=TRUE}, which internally calls the function \texttt{aggregateReference()}. This function performs k-means clustering to aggregate reference cells prior to annotation, and the initialization of k-means depends on randomly selected starting points. Since no fixed seed was specified before this procedure, independent executions generated slightly different pseudo-bulk reference profiles and, consequently, potentially different annotation labels. Although the use of \texttt{aggr.ref=TRUE} significantly improved computational efficiency, it introduced variability between executions. Nevertheless, the practical impact of this issue is likely limited because the same reference dataset was consistently employed throughout the study, and strongly expressed marker genes remained stable across analyses. However, clustering-related outputs generated throughout the pipeline cannot be considered fully reproducible in the absence of fixed random seeds. This limitation was identified during the final stages of the project, and rerunning the complete annotation workflow was not feasible within the available timeframe due to the substantial computational resources and execution times required.

The entirety of this thesis relies on the assumption that the samples were derived from MIBC tissue and that the associated metadata (covering treatment, response, and tumor presence) was accurate.
%%%%%%%%%%%%%%%%%%%%%%%%%%%%%%%%%%%%%%%%%%%%%%%%%%%%%%%%%%%%%%%%%%%%%%%%%%%%%%%%%%%%%%%%%%%%%%%

% In this chapter I will describe the most common options used, both the 
% ones inherited from \Class{scrbook} and the \Class{kao}-specific ones. 
% Options passed to the class modifies its default behaviour; beware 
% though that some options may lead to unexpected results\ldots

% \section{\Class{KOMA} Options}

% The \Class{kaobook} class is based on \Class{scrbook}, therefore it 
% understands all of the options you would normally pass to that class. If 
% you have a lot of patience, you can read the \KOMAScript\xspace 
% guide.\sidenote{The guide can be downloaded from 
% \url{https://ctan.org/pkg/koma-script?lang=en}.} Actually, the reading 
% of such guide is suggested as it is very instructive.

% Every \KOMAScript\xspace option you pass to the class when you load it 
% is automatically activated. In addition, in \Class{kaobook} some options 
% have modified default values. For instance, the font size is 9.5pt and 
% the paragraphs are separated by space,\sidenote[][-7mm]{To be precise, 
% they are separated by half a line worth of space: the \Option{parskip} 
% value is \enquote{half}.} not marked by indentation.

% \section{\Class{kao} Options}

% In the future I plan to add more options to set the paragraph formatting 
% (justified or ragged) and the position of the margins (inner or outer in 
% twoside mode, left or right in oneside mode).\sidenote{As of now, 
% paragraphs are justified, formatted with \Command{singlespacing} (from 
% the \Package{setspace} package) and \Command{frenchspacing}.}

% I take this opportunity to renew the call for help: everyone is 
% encouraged to add features or reimplement existing ones, and to send me 
% the results. You can find the GitHub repository at 
% \url{https://github.com/fmarotta/kaobook}.

% \begin{kaobox}[frametitle=To Do]
% Implement the \Option{justified} and \Option{margin} options. To be 
% consistent with the \KOMAScript\xspace style, they should accept a 
% simple switch as a parameter, where the simple switch should be 
% \Option{true} or \Option{false}, or one of the other standard values for 
% simple switches supported by \KOMAScript. See the \KOMAScript\xspace 
% documentation for further information.
% \end{kaobox}

% The above box is an example of a \Environment{kaobox}, which will be 
% discussed more thoroughly in \frefch{mathematics}. Throughout the book I 
% shall use these boxes to remarks what still needs to be done.

% \section{Other Things Worth Knowing}

% A bunch of packages are already loaded in the class because they are 
% needed for the implementation. These include:

% \begin{itemize}
% 	\item etoolbox
% 	\item calc
% 	\item xifthen
% 	\item xkeyval
% 	\item xparse
% 	\item xstring
% \end{itemize}

% Many more packages are loaded, but they will be discussed in due time. 
% Here, we will mention only one more set of packages, needed to change 
% the paragraph formatting (recall that in the future there will be 
% options to change this). In particular, the packages we load are:

% \begin{itemize}
% 	\item ragged2e
% 	\item setspace
% 	\item hyphenat
% 	\item microtype
% 	\item needspace
% 	\item xspace
% 	\item xcolor (with options \Option{usenames,dvipsnames})
% \end{itemize}

% Some of the above packages do not concern paragraph formatting, but we 
% nevertheless grouped them with the others. By default, the main text is 
% justified and formatted with singlespacing and frenchspacing; the margin 
% text is the same, except that the font is a bit smaller.

% As a last warning, please be aware that the \Package{cleveref} package 
% is not compatible with \Class{kaobook}. You should use the commands 
% discussed in \refsec{hyprefs} instead.

% \section{Document Structure}

% We provide optional arguments to the \Command{title} and 
% \Command{author} commands so that you can insert short, plain text 
% versions of this fields, which can be used, typically in the half-title 
% or somewhere else in the front matter, through the commands 
% \Command{@plaintitle} and \Command{@plainauthor}, respectively. The PDF 
% properties \Option{pdftitle} and \Option{pdfauthor} are automatically 
% set by hyperref to the plain values if present, otherwise to the normal 
% values.\sidenote[][*-1]{We think that this is an important point so 
% we remark it here. If you compile the document with pdflatex, the PDF 
% metadata will be altered so that they match the plain title and author 
% you have specified; if you did not specify them, the metadata will be 
% set to the normal title and author.}

% There are defined two page layouts, \Option{margin} and \Option{wide}, 
% and two page styles, \Option{plain} and \Option{fancy}. The layout 
% basically concern the width of the margins, while the style refers to 
% headers and footer; these issues will be 
% discussed in \frefch{layout}.\sidenote[][6mm]{For now, suffice it to say that pages with 
% the \Option{margin} layout have wide margins, while with the 
% \Option{wide} layout the margins are absent. In \Option{plain} pages the 
% headers and footer are suppressed, while in \Option{fancy} pages there 
% is a header.} 

% The commands \Command{frontmatter}, \Command{mainmatter}, and 
% \Command{backmatter} have been redefined in order to automatically 
% change page layout and style for these sections of the book. The front 
% matter uses the \Option{margin} layout and the \Option{plain} page 
% style. In the mainmatter the margins are wide and the headings are 
% fancy. In the appendix the style and the layout do not change; however 
% we use \Command{bookmarksetup\{startatroot\}} so that the bookmarks of 
% the chapters are on the root level (without this, they would be under 
% the preceding part). In the backmatter the margins shrink again and we 
% also reset the bookmarks root.
